\chapter{概率方法、随机舍入与随机游走}

本章展示随机性除“直接运行算法”外的三种作用：证明某个组合对象必然存在、把连续解转成离散解、在状态空间上采样。

\section{概率方法}

概率方法的基本逻辑是：若随机生成对象满足性质 $P$ 的概率大于 0，那么至少存在一个确定对象满足 $P$。证明可以完全非构造，也可以进一步导出随机算法或确定性算法。

\begin{example}[Max-Cut 的 $m/2$ 保证]
对无向图的每个顶点独立抛公平硬币，决定放入割的左侧或右侧。每条边跨越割的概率为 $1/2$。若 $X$ 是跨割边数，由期望线性性
\[
\E X=\frac m2.
\]
因此至少存在一个割包含不少于 $m/2$ 条边；否则所有割都少于 $m/2$，平均值也会小于 $m/2$。
\end{example}

\begin{example}[随机赋值满足 3-CNF 子句]
若一个子句包含三个互异文字，随机独立赋值使它不满足的概率为 $1/8$，满足概率为 $7/8$。因此对 $m$ 个这类子句，满足子句数的期望为 $7m/8$，必存在至少满足该数量的赋值。
\end{example}

\section{条件期望法}

概率方法常可去随机化。若目标随机变量为 $X$，逐个固定随机选择，并总是选择使条件期望不下降的分支：
\[
\E[X|R_1=r_1,\dots,R_t=r_t]
\ge
\E[X|R_1=r_1,\dots,R_{t-1}=r_{t-1}].
\]
因为父节点条件期望是各分支条件期望的加权平均，至少一个分支不低于父节点。全部固定后，实际目标值不低于初始期望。

对 Max-Cut，可以逐个决定顶点放哪侧，计算尚未决定边的条件期望，得到确定性 $m/2$ 保证。

\begin{intuitionbox}
“随机解的期望很好”不一定意味着必须真的抛硬币。若条件期望能高效计算，就可以沿着不降低期望的分支确定性地走到一个好解。
\end{intuitionbox}

\section{随机舍入}

组合优化常先求线性规划松弛，再把分数解随机转成 0/1 解。设集合覆盖实例中，元素宇宙有 $n$ 个元素，集合 $S_j$ 成本为 $c_j$，LP 解满足
\[
\sum_{j:e\in S_j}x_j\ge1,
\qquad 0\le x_j\le1.
\]
令 $\alpha=\ln n+1$，独立选择集合 $S_j$ 的概率为
\[
p_j=\min\{1,\alpha x_j\}.
\]
期望成本满足
\[
\E[\text{cost}]=\sum_jc_jp_j
\le\alpha\sum_jc_jx_j.
\]
对任意元素 $e$，若没有某个 $p_j=1$ 的覆盖集合，则它未被覆盖的概率
\begin{align*}
\Prb[e\text{ 未覆盖}]
&=\prod_{j:e\in S_j}(1-p_j)\\
&\le\exp\left(-\sum_{j:e\in S_j}p_j\right)\\
&\le e^{-\alpha}<\frac1n.
\end{align*}
再用 union bound 控制存在未覆盖元素的概率。实际算法可以重试，或对少量未覆盖元素做确定性修补。

\begin{warningbox}
随机舍入的独立选择、LP 约束、缩放因子和修补步骤都是保证的一部分。只写“按 $x_j$ 当概率采样”通常不足以得到可行解或高概率界。
\end{warningbox}

\section{Markov 链基础}

有限状态 Markov 链由转移矩阵 $P$ 定义：
\[
P(x,y)=\Prb[X_{t+1}=y|X_t=x],
\qquad
\sum_yP(x,y)=1.
\]
Markov 性表示给定当前状态后，下一步与更早历史条件独立。

\begin{definition}[平稳分布]
分布 $\pi$ 若满足
\[
\pi(y)=\sum_x\pi(x)P(x,y),
\quad\text{即}\quad \pi=\pi P,
\]
则称为平稳分布。
\end{definition}

若链不可约且非周期，则从任意初始分布出发会收敛到唯一平稳分布。实际常加入“以概率 $1/2$ 原地不动”的 lazy 步骤消除二分图随机游走的周期性。

\section{无向图随机游走}

在无向图 $G=(V,E)$ 上，从顶点 $u$ 均匀走到邻居：
\[
P(u,v)=
\begin{cases}
1/\deg(u),&(u,v)\in E,\\
0,&\text{否则}.
\end{cases}
\]
若图连通，候选平稳分布为
\[
\pi(v)=\frac{\deg(v)}{2|E|}.
\]
验证 detailed balance：对边 $(u,v)$，
\[
\pi(u)P(u,v)=\frac1{2|E|}=\pi(v)P(v,u).
\]
对所有 $u$ 求和即得 $\pi P=\pi$。

这说明普通随机游走更常访问高度顶点，不是对顶点均匀采样。若研究需要均匀分布，必须修改转移规则，例如使用 Metropolis--Hastings 校正。

\section{混合时间与 MCMC 边界}

总变差距离为
\[
\|\mu-\pi\|_{\mathrm{TV}}
=\frac12\sum_x|\mu(x)-\pi(x)|.
\]
混合时间 $t_{\mathrm{mix}}(\epsilon)$ 是从最坏初态出发，使分布距离平稳分布至多 $\epsilon$ 所需的最小步数。

MCMC 用容易模拟的 Markov 链近似从复杂分布采样。仅验证平稳分布还不够：链可能混合极慢，有限步样本仍强烈依赖初始状态。实验必须记录 burn-in、链长度、自相关和多链诊断，并且不能用“轨迹看起来稳定”替代理论混合界。

\section{小结与验收}

\begin{checkpointbox}
\begin{enumerate}
  \item 用期望线性性证明存在大小至少 $m/2$ 的割；
  \item 说明条件期望法为什么每一步总能选到不变差的分支；
  \item 在集合覆盖随机舍入中推导单个元素未覆盖概率；
  \item 验证无向图随机游走的平稳分布；
  \item 给出一个“平稳分布正确但有限步采样仍不可信”的原因。
\end{enumerate}
\end{checkpointbox}
